% TOP_PROBLEM: {"id": 30007286, "title": "The two-primary BSD formula at low two-Selmer corank", "category_id": 1, "category": {"id": 1, "name": "number_theory", "display_name": "Number Theory", "description": "Properties of integers, prime numbers, Diophantine equations.", "slug": "number-theory", "order_index": 1, "created_at": "2026-07-31T15:26:25.670Z"}, "status": "open"}
% FIRST_POSTED: 2026-10-06
% ATTEMPT_STATUS: solved
% OPENAI_CLAIM: {"schema_version":1,"source_repo":"https://github.com/openai/math","source_commit":"adc7f1241b42e322a6451854ab7e4b4c146bf78a","release_date":"2026-10-06","match":"full","resolution":"proved","independently_reviewed":false,"second_pass":"agreed","families":[{"family":"002","title":"The full BSD formula from low Selmer corank","subject":"Number theory","manuscripts":[{"title":"The two-primary Birch–Swinnerton-Dyer formula in Selmer corank at most one","pdf_path":"preprints/The-two-primary-Birch-Swinnerton-Dyer-formula-in-Selmer-corank-at-most-one-September-24-2026/paper.pdf","readme_path":"preprints/The-two-primary-Birch-Swinnerton-Dyer-formula-in-Selmer-corank-at-most-one-September-24-2026/README.md","date":"2026-09-24","theorem_ref":"thm:main (sections/01-introduction.tex:54)","statement_sources":[{"path":"preprints/The-two-primary-Birch-Swinnerton-Dyer-formula-in-Selmer-corank-at-most-one-September-24-2026/build/sections/01-introduction.tex","line":54,"label":"thm:main"}]}],"lean":null,"reasoning_trace":null}]}
\documentclass[11pt]{article}
\usepackage[margin=25mm]{geometry}
\usepackage{fontspec}
\setmainfont{Latin Modern Roman}
\usepackage{amsmath,amssymb,mathtools}
\usepackage{unicode-math}
\setmathfont{Latin Modern Math}
\usepackage{enumitem}
\usepackage{xurl,hyperref,bookmark}
\hypersetup{colorlinks=true,urlcolor=blue}
\setcounter{secnumdepth}{0}
\setlength{\parindent}{0pt}
\setlength{\parskip}{5pt}
\setlength{\emergencystretch}{3em}
\begin{document}
\section*{The two-primary BSD formula at low two-Selmer corank}

\subsection{OpenAI's claim}
OpenAI presents the following claims relevant to this question, with manuscript numbers following paper order in Sources. In manuscript 1 (family 002; 2026-09-24), the asserted result is that the two-adic part of the BSD leading-coefficient identity holds for rational elliptic curves of two-power Selmer corank at most one, and for almost every quadratic twist.

\subsection{Mathematical statement of the claimed result}
\paragraph{Manuscript 1.}
For any elliptic curve $E/\mathbb Q$, let $s_2(E)=\operatorname{corank}_{\mathbb Z_2}\operatorname{Sel}_{2^{\infty}}(E/\mathbb Q)$ and assume $s_2(E)\le1$. The claim is
\[
 r=\operatorname{rank}E(\mathbb Q)=\operatorname{ord}_{s=1}L(E,s)=s_2(E),
 \qquad |\operatorname{Sha}(E/\mathbb Q)|<\infty.
\]
With $T_E=E(\mathbb Q)_{\rm tors}$, real N\'eron period $\Omega_E$, N\'eron--Tate regulator $\operatorname{Reg}_E$ and finite-prime Tamagawa numbers $c_{\ell}(E)$, it further claims
\[
 Q_E=\frac{L^{(r)}(E,1)|T_E|^2}{r!\Omega_E\operatorname{Reg}_E\prod_{\ell} c_{\ell}(E)}
 \in\mathbb Q_{>0},\qquad v_2(Q_E)=v_2(|\operatorname{Sha}(E/\mathbb Q)|).
\]
There is no additional restriction on $E[2]$ or reduction at two. The formula determines the two-primary valuation; it does not state equality $Q_E=|\operatorname{Sha}|$ at all primes or address larger Selmer corank.

\subsection{Outline of the argument}
For manuscript 1 (family 002; 2026-09-24), the proof compares integral determinant lattices with Siegel-unit periods and central Kummer lattices. Separate residual and character-division arguments establish primitivity at two, and the rank-one comparison uses a squared Heegner logarithm rather than only rational period identities.

\subsection{Scope relative to this problem}
The result supplies the conclusion in this problem with its displayed hypotheses and quantifiers. The broader statements and applications in the manuscript are kept separate from the particular question above. The supporting result is manuscript 1 (family 002; 2026-09-24).

\subsection{Formal verification}
No Lean formalization is listed at this snapshot for manuscript 1 (family 002; 2026-09-24).

This site has not executed the Lean code. Described statement coverage is not an independent proof check.

\subsection{Status caveat}
The manuscripts were produced by an internal OpenAI model. OpenAI cautions that results without formalization can contain mistakes. This site has not independently checked the proofs; the label records the scope of OpenAI's claim.

\subsection{Sources}
\begin{itemize}
\item OpenAI, \emph{The two-primary Birch–Swinnerton-Dyer formula in Selmer corank at most one}, 2026-09-24; citation key \texttt{OAI:The-two-primary-Birch-Swinnerton-Dyer-formula-in-Selmer-corank-at-most-one-September-24-2026}. \href{https://github.com/openai/math/blob/adc7f1241b42e322a6451854ab7e4b4c146bf78a/preprints/The-two-primary-Birch-Swinnerton-Dyer-formula-in-Selmer-corank-at-most-one-September-24-2026/paper.pdf}{PDF (GitHub)}; \href{https://github.com/openai/math/raw/adc7f1241b42e322a6451854ab7e4b4c146bf78a/preprints/The-two-primary-Birch-Swinnerton-Dyer-formula-in-Selmer-corank-at-most-one-September-24-2026/paper.pdf}{PDF (direct)}; \href{https://github.com/openai/math/blob/adc7f1241b42e322a6451854ab7e4b4c146bf78a/preprints/The-two-primary-Birch-Swinnerton-Dyer-formula-in-Selmer-corank-at-most-one-September-24-2026/README.md}{Manuscript page and citation}; \href{https://github.com/openai/math/blob/main/preprints/The-two-primary-Birch-Swinnerton-Dyer-formula-in-Selmer-corank-at-most-one-September-24-2026/paper.pdf}{Latest version}; Statement source: \href{https://github.com/openai/math/blob/adc7f1241b42e322a6451854ab7e4b4c146bf78a/preprints/The-two-primary-Birch-Swinnerton-Dyer-formula-in-Selmer-corank-at-most-one-September-24-2026/build/sections/01-introduction.tex\#L54}{thm:main}.
\item \href{https://github.com/openai/math/blob/adc7f1241b42e322a6451854ab7e4b4c146bf78a/CONTENTS.md}{OpenAI catalogue entry}.
\end{itemize}
\end{document}
